% TOP_PROBLEM: {"id":30007562,"problem_number":"LOCAL-30007562","title":"Speed of learning in games","category_id":21,"category":{"id":21,"name":"economics","display_name":"Economics","description":"Open economic theory statements, published research agendas, and proposed scoped specifications; question provenance and historical priority are qualified per record.","slug":"economics","order_index":21,"created_at":"2026-10-08T00:00:00Z"},"status":"research_question","external_url":null,"legacy_ids":[],"published":false,"statement":"Estimate how long heterogeneous human players take to approach low exploitability, and predict that time across games and feedback regimes without assuming universal Nash convergence.","economics":{"original_id":51,"field":"Game theory and strategic interaction","kind":"empirical","formalization_basis":"adapted_research_question","source_status":"research_agenda","priority_status":"earliest_found","priority_note":"This is the earliest source in which the relevant agenda was verified during this audit; absolute priority has not been established.","earliest_verified_reference":{"authors":["Drew Fudenberg","David K. Levine"],"title":"Whither Game Theory? Towards a Theory of Learning in Games","year":2016,"url":"https://economics.mit.edu/sites/default/files/2022-09/Whither%20Game%20Theory%20Towards%20a%20Theory%20of%20Learning%20in%20Games.pdf","doi":"10.1257/jep.30.4.151","locator":"Introduction, p.153; conclusion, pp.165–166","evidence_summary":"The essay explicitly calls for more knowledge about learning speed, equilibrium prediction, and tractable models that allow reevaluation."},"current_reference":{"authors":["Drew Fudenberg","David K. Levine"],"title":"Whither Game Theory? Towards a Theory of Learning in Games","year":2016,"url":"https://economics.mit.edu/sites/default/files/2022-09/Whither%20Game%20Theory%20Towards%20a%20Theory%20of%20Learning%20in%20Games.pdf","doi":"10.1257/jep.30.4.151","locator":"Introduction, p.153; conclusion, pp.165–166","evidence_summary":"The essay explicitly calls for more knowledge about learning speed, equilibrium prediction, and tractable models that allow reevaluation."},"formalization_note":"The horizon functional, protocol, and out-of-game validation are an empirical adaptation of the explicit learning-speed agenda."},"statusQualification":"Published question or research agenda verified at the cited locator. The mathematical specification is adapted for this catalogue and includes additional assumptions; source publication does not establish first-ever priority or certify this exact model remains unsolved. The horizon functional, protocol, and out-of-game validation are an empirical adaptation of the explicit learning-speed agenda.","statusReviewedAt":"2026-10-08"}
% FIRST_POSTED: 2026-10-08
% ENTRY_KIND: statement_only
% !TeX program = lualatex
\documentclass[11pt]{article}
\usepackage{fontspec}
\usepackage[a4paper,margin=25mm]{geometry}
\usepackage{amsmath,amssymb,mathtools}
\usepackage{xurl}
\usepackage[hidelinks]{hyperref}
\newcommand{\sref}[2]{\hyperlink{source:#1:#2}{[S#2]}}
\setlength{\emergencystretch}{3em}
\begin{document}
% problemId:30007562
\section*{Speed of learning in games}\label{30007562}
\subsection{Definitions and mathematical statement}
Fix a class \(\mathcal G\) of finite normal-form games with payoffs \(u_i\in[0,1]\), an information regime \(f\), and a population \(P\) of human players. For game \(G\), repeatedly rematch or retain players according to a preregistered protocol, observe actions \(a_{it}\), and elicit beliefs when feasible. Let \(\widehat\sigma_T\) be the product of each player's empirical action distribution through period \(T\). Define exploitability
\[
 e_G(\sigma)=\max_i\max_{b_i}\{u_i(b_i,\sigma_{-i})-u_i(\sigma)\},
\]
where \(b_i\) ranges over pure actions and payoffs to mixed profiles are expected payoffs. For \(\varepsilon>0\) and \(\delta\in(0,1)\), define \(T_{\varepsilon,\delta}(G,f,P)\) as the least horizon after which the probability of \(e_G(\widehat\sigma_t)>\varepsilon\) is at most \(\delta\) for every measured \(t\). Estimate this horizon or establish that the experimental observation window only gives a lower bound. Determine how feedback, game structure, and population heterogeneity affect it, and test models trained on different games. Dependence across periods and multiple horizon searches require valid uncertainty accounting. Product-marginal exploitability must be distinguished from the regret of the empirical joint distribution, which is another observable. The agenda is finite-horizon human prediction; it does not assume that human play in every game converges to Nash equilibrium.

\par\textbf{Formulation and scope.} The horizon functional, protocol, and out-of-game validation are an empirical adaptation of the explicit learning-speed agenda.

\par\textbf{Open-question provenance.} An explicit question or research agenda was verified at the source locator. This does not by itself certify that every later version remains unresolved \sref{30007562}{1}.

\par\textbf{Historical priority.} This is the earliest source in which the relevant agenda was verified during this audit; absolute priority has not been established.

\subsection{Short English statement}
Estimate how long heterogeneous human players take to approach low exploitability, and predict that time across games and feedback regimes without assuming universal Nash convergence.

\subsection{Sources}
\begin{itemize}\raggedright
\item[\textnormal{[S1]}]\hypertarget{source:30007562:1}{} Drew Fudenberg, David K. Levine. 2016. Whither Game Theory? Towards a Theory of Learning in Games. Introduction, p.153; conclusion, pp.165–166.
\url{https://economics.mit.edu/sites/default/files/2022-09/Whither%20Game%20Theory%20Towards%20a%20Theory%20of%20Learning%20in%20Games.pdf}.
DOI: 10.1257/jep.30.4.151.
{\small\textit{Earliest verified question/agenda reference; historical priority qualified below. The essay explicitly calls for more knowledge about learning speed, equilibrium prediction, and tractable models that allow reevaluation.}}
\end{itemize}
\end{document}
